\chapter{مقدمه}

برنامه‌نویسی رقابتی ترکیب دو مبحث است:
(۱) طراحی الگوریتم‌ها و (۲) پیاده‌سازی الگوریتم‌ها.

\key{طراحی الگوریتم‌ها} شامل حل مسئله
و تفکر ریاضی است.
مهارت‌های تحلیل مسائل و حل خلاقانه آن‌ها
مورد نیاز است.
الگوریتم حل مسئله
باید هم درست و هم کارآمد باشد،
و هسته مسئله اغلب
ابداع الگوریتمی کارآمد است.

دانش نظری الگوریتم‌ها
برای برنامه‌نویسان رقابتی مهم است.
معمولا راه‌حل مسئله
ترکیبی از روش‌های شناخته‌شده و
بینش‌های جدید است.
تکنیک‌های حاضر در برنامه‌نویسی رقابتی
اساس تحقیقات علمی
الگوریتم‌ها را نیز تشکیل می‌دهند.

\key{پیاده‌سازی الگوریتم‌ها} نیازمند مهارت‌های برنامه‌نویسی
خوب است.
در برنامه‌نویسی رقابتی، راه‌حل‌ها
با ارزیابی الگوریتم پیاده‌سازی‌شده
توسط مجموعه‌ای از تست‌ها داوری می‌شوند.
بنابراین، تنها درست بودن ایده الگوریتم
کافی نیست، بلکه پیاده‌سازی نیز
باید درست باشد.

سبک کدنویسی خوب در مسابقات
سرراست و مختصر است.
برنامه‌ها باید سریع نوشته شوند،
زیرا زمان زیادی وجود ندارد.
برخلاف مهندسی نرم‌افزار سنتی،
برنامه‌ها کوتاه‌اند (معمولا حداکثر چند
صد خط کد)، و پس از مسابقه نیازی به
نگهداری آن‌ها نیست.

\section{زبان‌های برنامه‌نویسی}

\index{programming language@زبان برنامه‌نویسی}

در حال حاضر، محبوب‌ترین زبان‌های
برنامه‌نویسی در مسابقات C++، پایتون و جاوا هستند.
برای نمونه، در Google Code Jam 2017،
میان ۳۰۰۰ شرکت‌کننده برتر،
۷۹ ٪ از C++،
۱۶ ٪ از پایتون و
۸ ٪ از جاوا استفاده کردند \cite{goo17}.
برخی شرکت‌کنندگان از چند زبان نیز استفاده کردند.

بسیاری بر این باورند که C++ بهترین انتخاب
برای برنامه‌نویس رقابتی است،
و C++ تقریبا همیشه در
سامانه‌های مسابقه در دسترس است.
مزایای استفاده از C++ این است که
زبانی بسیار کارآمد است و
کتابخانه استاندارد آن شامل مجموعه‌ای
بزرگ از داده‌ساختارها و الگوریتم‌ها است.

از سوی دیگر، تسلط بر
چندین زبان و درک
نقاط قوت آن‌ها مفید است.
برای مثال، اگر اعداد صحیح بزرگ در
مسئله نیاز باشند،
پایتون می‌تواند انتخاب خوبی باشد، زیرا
عملیات توکار برای
محاسبات با اعداد صحیح بزرگ دارد.
با این حال، بیشتر مسائل در مسابقات برنامه‌نویسی
طوری طرح می‌شوند که
استفاده از زبانی خاص، مزیت ناعادلانه نباشد.

تمام برنامه‌های نمونه در این کتاب به C++ نوشته شده‌اند،
و اغلب از داده‌ساختارها و الگوریتم‌های کتابخانه
استاندارد استفاده می‌شود.
برنامه‌ها از استاندارد C++11 پیروی می‌کنند،
که امروزه در اکثر مسابقات قابل استفاده است.
اگر هنوز نمی‌توانید با C++ برنامه‌نویسی کنید،
اکنون زمان مناسبی برای یادگیری آن است.

\subsubsection{قالب کد C++}

قالب کد معمول C++ برای برنامه‌نویسی رقابتی
به این شکل است:

\begin{lstlisting}
#include <bits/stdc++.h>

using namespace std;

int main() {
    // solution comes here
}
\end{lstlisting}

خط \texttt{\#include} در ابتدای کد قابلیتی از کامپایلر \texttt{g++} است که اجازه می‌دهد کل کتابخانه استاندارد اضافه شود. بنابراین نیازی به اضافه کردن جداگانه کتابخانه‌هایی مانند \texttt{iostream}، \texttt{vector} و \texttt{algorithm} نیست، بلکه خودکار در دسترس قرار می‌گیرند.

خط \texttt{using} اعلام می‌کند که کلاس‌ها و توابع کتابخانه استاندارد می‌توانند مستقیماً در کد استفاده شوند. بدون خط \texttt{using} باید مثلاً \texttt{std::cout} نوشته می‌شد، اما اکنون نوشتن \texttt{cout} کفایت می‌کند.

کد را می‌توان با استفاده از دستور زیر کامپایل کرد:

\begin{lstlisting}
g++ -std=c++11 -O2 -Wall test.cpp -o test
\end{lstlisting}

این دستور فایل دودویی \texttt{test} را از کد منبع \texttt{test.cpp} تولید می‌کند. کامپایلر از استاندارد C++11 پیروی می‌کند (\texttt{-std=c++11})، کد را بهینه‌سازی می‌کند (\texttt{-O2}) و هشدارهایی درباره خطاهای احتمالی نمایش می‌دهد (\texttt{-Wall}).

\section{ورودی و خروجی}

\index{ورودی و خروجی}

در بیشتر مسابقات، جریان‌های استاندارد برای خواندن ورودی و نوشتن خروجی استفاده می‌شوند. در C++ جریان‌های استاندارد عبارتند از \texttt{cin} برای ورودی و \texttt{cout} برای خروجی. علاوه بر این، توابع C یعنی \texttt{scanf} و \texttt{printf} نیز می‌توانند به کار روند.

ورودی برنامه معمولاً شامل اعداد و رشته‌هایی است که با فاصله و خطوط جدید جدا شده‌اند. آن‌ها را می‌توان از جریان \texttt{cin} به صورت زیر خواند:

\begin{lstlisting}
int a, b;
string x;
cin >> a >> b >> x;
\end{lstlisting}

این نوع کد همیشه کار می‌کند، به این فرض که حداقل یک فاصله یا خط جدید بین هر عنصر در ورودی وجود داشته باشد. برای مثال، کد بالا هر دو ورودی زیر را می‌خواند:
\begin{lstlisting}
123 456 monkey
\end{lstlisting}
\begin{lstlisting}
123    456
monkey
\end{lstlisting}
جریان \texttt{cout} برای خروجی به شکل زیر استفاده می‌شود:
\begin{lstlisting}
int a = 123, b = 456;
string x = "monkey";
cout << a << " " << b << " " << x << "\n";
\end{lstlisting}

ورودی و خروجی گاهی گلوگاه برنامه است. خطوط زیر در ابتدای کد، ورودی و خروجی را کارآمدتر می‌کنند:

\begin{lstlisting}
ios::sync_with_stdio(0);
cin.tie(0);
\end{lstlisting}

توجه کنید که کاراکتر خط جدید \texttt{"\textbackslash n"} سریع‌تر از \texttt{endl} عمل می‌کند، چون \texttt{endl} همیشه یک عملیات فلاش (تخلیه بافر) انجام می‌دهد.

توابع C یعنی \texttt{scanf} و \texttt{printf} جایگزینی برای جریان‌های استاندارد C++ هستند. آن‌ها معمولاً کمی سریع‌ترند، ولی استفاده از آن‌ها سخت‌تر است. کد زیر دو عدد صحیح را از ورودی می‌خواند:
\begin{lstlisting}
int a, b;
scanf("%d %d", &a, &b);
\end{lstlisting}
کد زیر دو عدد صحیح را چاپ می‌کند:
\begin{lstlisting}
int a = 123, b = 456;
printf("%d %d\n", a, b);
\end{lstlisting}

گاهی برنامه باید یک خط کامل را از ورودی بخواند که ممکن است شامل فاصله باشد. این کار با استفاده از تابع \texttt{getline} انجام می‌شود:

\begin{lstlisting}
string s;
getline(cin, s);
\end{lstlisting}

اگر مقدار داده‌ها نامشخص باشد، حلقه زیر مفید است:
\begin{lstlisting}
while (cin >> x) {
    // code
}
\end{lstlisting}
این حلقه عناصر را یکی پس از دیگری از ورودی می‌خواند، تا زمانی که داده دیگری در ورودی وجود نداشته باشد.

در برخی سامانه‌های مسابقه، از فایل‌ها برای ورودی و خروجی استفاده می‌شود.
راه حل آسان برای این موضوع، نوشتن کد به شکل همیشگی با جریان‌های استاندارد و اضافه کردن خطوط زیر به ابتدای کد است:
\begin{lstlisting}
freopen("input.txt", "r", stdin);
freopen("output.txt", "w", stdout);
\end{lstlisting}
پس از آن، برنامه ورودی را از فایل ''input.txt'' خوانده و خروجی را در فایل ''output.txt'' می‌نویسد.

\section{کار با اعداد}

\index{integer}

\subsubsection{اعداد صحیح}

پرکاربردترین نوع عدد صحیح در برنامه‌نویسی رقابتی \texttt{int} است؛ نوعی ۳۲ بیتی با محدوده مقادیر $-2^{31} \ldots 2^{31}-1$ یا حدود $-2 \cdot 10^9 \ldots 2 \cdot 10^9$.
اگر نوع \texttt{int} کافی نباشد، می‌توان از نوع ۶۴ بیتی \texttt{long long} استفاده کرد.
این نوع دارای محدوده مقادیر $-2^{63} \ldots 2^{63}-1$ یا حدود $-9 \cdot 10^{18} \ldots 9 \cdot 10^{18}$ است.

کد زیر یک متغیر از نوع \texttt{long long} تعریف می‌کند:
\begin{lstlisting}
long long x = 123456789123456789LL;
\end{lstlisting}
پسوند \texttt{LL} نشان می‌دهد که نوع عدد \texttt{long long} است.

یک اشتباه متداول هنگام استفاده از نوع \texttt{long long} این است که نوع \texttt{int} همچنان در بخش دیگری از کد استفاده شود.
برای نمونه، کد زیر حاوی یک خطای ظریف است:

\begin{lstlisting}
int a = 123456789;
long long b = a*a;
cout << b << "\n"; // -1757895751
\end{lstlisting}

اگرچه متغیر \texttt{b} از نوع \texttt{long long} است، اما هر دو عدد در عبارت \texttt{a*a} از نوع \texttt{int} هستند و نتیجه نیز از نوع \texttt{int} خواهد بود.
به همین دلیل، متغیر \texttt{b} نتیجه نادرستی ذخیره می‌کند.
مشکل با تغییر نوع \texttt{a} به \texttt{long long} یا تغییر عبارت به \texttt{(long long)a*a} حل می‌شود.

معمولاً مسائل مسابقات به گونه‌ای طراحی می‌شوند که نوع \texttt{long long} کافی باشد.
با این حال، مفید است بدانید کامپایلر \texttt{g++} نوع ۱۲۸ بیتی \texttt{\_\_int128\_t} را نیز با محدوده مقادیر $-2^{127} \ldots 2^{127}-1$ یا حدود $-10^{38} \ldots 10^{38}$ ارائه می‌دهد.
البته این نوع در تمام سامانه‌های مسابقه در دسترس نیست.

\subsubsection{حساب پیمانه‌ای}

\index{remainder}
\index{modular arithmetic}

باقی‌مانده تقسیم $x$ بر $m$ با $x \bmod m$ نشان داده می‌شود.
برای مثال، $17 \bmod 5 = 2$، زیرا $17 = 3 \cdot 5 + 2$.

گاهی پاسخ مسئله عددی بسیار بزرگ است اما چاپ آن «به پیمانه $m$» کفایت می‌کند، یعنی باقی‌مانده تقسیم پاسخ بر $m$ (برای نمونه، «به پیمانه $10^9+7$»).
ایده این است که حتی در صورت بسیار بزرگ بودن پاسخ نهایی، استفاده از نوع‌های \texttt{int} و \texttt{long long} کافی باشد.

یک ویژگی مهم باقی‌مانده این است که
در جمع، تفریق و ضرب،
می‌توان باقی‌مانده را قبل از انجام عملیات محاسبه کرد:

\[
\begin{array}{rcr}
(a+b) \bmod m & = & (a \bmod m + b \bmod m) \bmod m \\
(a-b) \bmod m & = & (a \bmod m - b \bmod m) \bmod m \\
(a \cdot b) \bmod m & = & (a \bmod m \cdot b \bmod m) \bmod m
\end{array}
\]

بنابراین، می‌توان پس از هر عملیات باقی‌مانده را محاسبه کرد
تا اعداد هرگز بیش از حد بزرگ نشوند.

برای نمونه، کد زیر $n!$،
یعنی فاکتوریل $n$، را به پیمانه $m$ محاسبه می‌کند:
\begin{lstlisting}
long long x = 1;
for (int i = 2; i <= n; i++) {
    x = (x*i)%m;
}
cout << x%m << "\n";
\end{lstlisting}

معمولاً می‌خواهیم باقی‌مانده همواره
بین $0\ldots m-1$ باشد.
اما در C++ و زبان‌های دیگر،
باقی‌مانده یک عدد منفی
یا صفر است یا منفی.
یک روش آسان برای اطمینان از عدم وجود
باقی‌مانده‌های منفی این است که ابتدا باقی‌مانده را
به صورت معمول حساب کرده و در صورت منفی بودن نتیجه،
$m$ را به آن اضافه کنیم:
\begin{lstlisting}
x = x%m;
if (x < 0) x += m;
\end{lstlisting}
با این حال، این کار تنها زمانی لازم است که
در کد عملیات تفریق وجود داشته باشد و
احتمال منفی شدن باقی‌مانده برود.

\subsubsection{اعداد ممیز شناور}

\index{floating point number}

نوع‌های داده ممیز شناور معمول در
برنامه‌نویسی رقابتی عبارتند از
نوع ۶۴ بیتی \texttt{double}
و، به عنوان یک افزونه در کامپایلر \texttt{g++}،
نوع ۸۰ بیتی \texttt{long double}.
در بیشتر موارد، \texttt{double} کافی است،
اما \texttt{long double} دقت بالاتری دارد.

دقت مورد نیاز برای پاسخ معمولاً
در صورت مسئله ذکر می‌شود.
یک روش آسان برای چاپ پاسخ، استفاده از
تابع \texttt{printf}
و تعیین تعداد ارقام اعشار
در رشته قالب‌بندی است.
برای نمونه، کد زیر مقدار $x$ را
با ۹ رقم اعشار چاپ می‌کند:

\begin{lstlisting}
printf("%.9f\n", x);
\end{lstlisting}

یکی از دشواری‌های کار با اعداد ممیز شناور
این است که برخی اعداد را نمی‌توان به طور دقیق
به صورت ممیز شناور نمایش داد،
و خطاهای گرد کردن رخ خواهد داد.
برای نمونه، نتیجه کد زیر
غافلگیرکننده است:

\begin{lstlisting}
double x = 0.3*3+0.1;
printf("%.20f\n", x); // 0.99999999999999988898
\end{lstlisting}

به دلیل خطای گرد کردن،
مقدار \texttt{x} کمی کوچک‌تر از ۱ است،
در حالی که مقدار درست باید ۱ باشد.

مقایسه اعداد ممیز شناور با عملگر \texttt{==}
پرخطر است،
زیرا ممکن است مقادیر باید برابر باشند
اما به دلیل خطاهای دقت، برابر نباشند.
روش بهتر برای مقایسه اعداد ممیز شناور
این است که فرض کنیم دو عدد برابرند
اگر اختلاف بین آن‌ها کمتر از $\varepsilon$ باشد،
که در آن $\varepsilon$ یک عدد کوچک است.

در عمل، اعداد را می‌توان به شکل زیر
مقایسه کرد ($\varepsilon=10^{-9}$):

\begin{lstlisting}
if (abs(a-b) < 1e-9) {
    // a and b are equal
}
\end{lstlisting}

توجه کنید با این که اعداد ممیز شناور نادقیق هستند،
اعداد صحیح تا حد مشخصی دقیق ذخیره می‌شوند.
مثال: نوع \texttt{double} تمام اعداد صحیح
با قدر مطلق حداکثر $2^{53}$ را دقیق نمایش می‌دهد.

\section{کوتاه‌نویسی کد}

کد کوتاه در برنامه‌نویسی رقابتی ایده‌آل است،
چون کدها باید سریع نوشته شوند.
به همین دلیل، برنامه‌نویسان نام‌های کوتاه‌تری
برای انواع داده و اجزای کد تعریف می‌کنند.

\subsubsection{نام‌های نوع داده}
\index{tuppdef@\texttt{typedef}}
دستور \texttt{typedef} نام کوتاه‌تر
برای یک نوع داده می‌سازد.
مثال: نام \texttt{long long} طولانی است؛
نام کوتاه‌تر \texttt{ll} را تعریف می‌کنیم:
\begin{lstlisting}
typedef long long ll;
\end{lstlisting}
بعد از این، کد:
\begin{lstlisting}
long long a = 123456789;
long long b = 987654321;
cout << a*b << "\n";
\end{lstlisting}
این‌گونه کوتاه می‌شود:
\begin{lstlisting}
ll a = 123456789;
ll b = 987654321;
cout << a*b << "\n";
\end{lstlisting}

دستور \texttt{typedef} برای انواع داده پیچیده‌تر نیز کار می‌کند.
کد زیر نام \texttt{vi} را برای بردای از اعداد صحیح
و نام \texttt{pi} را برای زوج مرتب دو عدد صحیح تعریف می‌کند:
\begin{lstlisting}
typedef vector<int> vi;
typedef pair<int,int> pi;
\end{lstlisting}

\subsubsection{ماکروها}
\index{macro}
روش دیگر کوتاه‌نویسی، تعریف \key{ماکروها} است.
ماکرو رشته‌های مشخصی از کد را
قبل از کامپایل جایگزین می‌کند.
در C++، ماکروها با کلمه کلیدی \texttt{\#define} تعریف می‌شوند.

مثال: تعریف ماکروهای زیر:
\begin{lstlisting}
#define F first
#define S second
#define PB push_back
#define MP make_pair
\end{lstlisting}
بعد از این، کد:
\begin{lstlisting}
v.push_back(make_pair(y1,x1));
v.push_back(make_pair(y2,x2));
int d = v[i].first+v[i].second;
\end{lstlisting}
این‌گونه کوتاه می‌شود:
\begin{lstlisting}
v.PB(MP(y1,x1));
v.PB(MP(y2,x2));
int d = v[i].F+v[i].S;
\end{lstlisting}

ماکرو پارامتر هم می‌گیرد
که حلقه‌ها و ساختارها را کوتاه می‌کند.
مثال: تعریف ماکروی زیر:
\begin{lstlisting}
#define REP(i,a,b) for (int i = a; i <= b; i++)
\end{lstlisting}
بعد از این، کد:
\begin{lstlisting}
for (int i = 1; i <= n; i++) {
    search(i);
}
\end{lstlisting}
این‌گونه کوتاه می‌شود:
\begin{lstlisting}
REP(i,1,n) {
    search(i);
}
\end{lstlisting}

گاهی ماکروها خطاهایی می‌سازند که پیدا کردنشان سخت است.
مثال: ماکروی زیر برای محاسبه مربع یک عدد:
\begin{lstlisting}
#define SQ(a) a*a
\end{lstlisting}
این ماکرو همیشه درست کار \emph{نمی‌کند}.
مثال: کد
\begin{lstlisting}
cout << SQ(3+3) << "\n";
\end{lstlisting}
معادل این کد است:
\begin{lstlisting}
cout << 3+3*3+3 << "\n"; // 15
\end{lstlisting}

نسخه بهتر ماکرو به این صورت است:
\begin{lstlisting}
#define SQ(a) (a)*(a)
\end{lstlisting}
اکنون کد
\begin{lstlisting}
cout << SQ(3+3) << "\n";
\end{lstlisting}
معادل این کد است:
\begin{lstlisting}
cout << (3+3)*(3+3) << "\n"; // 36
\end{lstlisting}


\section{ریاضیات}

ریاضیات نقش مهمی در برنامه‌نویسی رقابتی دارد. موفقیت در برنامه‌نویسی رقابتی بدون مهارت‌های قوی ریاضی ممکن نیست. این بخش به مفاهیم و فرمول‌های مهم ریاضی می‌پردازد که در ادامه کتاب استفاده می‌شوند.

\subsubsection{فرمول‌های مجموع}

هر مجموع به فرم
\[\sum_{x=1}^n x^k = 1^k+2^k+3^k+\ldots+n^k,\]
که در آن $k$ عدد صحیح مثبت است، فرمول بسته دارد که چندجمله‌ای با درجه $k+1$ است.
برای مثال\footnote{\index{فرمول فالهابر}\index{Faulhaber's formula}
فرمول کلی برای این مجموع‌ها با نام \key{فرمول فالهابر} وجود دارد، ولی ارائه آن در اینجا پیچیده است.}،
\[\sum_{x=1}^n x = 1+2+3+\ldots+n = \frac{n(n+1)}{2}\]
و
\[\sum_{x=1}^n x^2 = 1^2+2^2+3^2+\ldots+n^2 = \frac{n(n+1)(2n+1)}{6}.\]

یک \key{تصاعد حسابی} \index{تصاعد حسابی}\index{arithmetic progression}
دنباله‌ای از اعداد است
که تفاضل هر دو عدد متوالی در آن ثابت است.
برای مثال،
\[3, 7, 11, 15\]
یک تصاعد حسابی با قدر نسبت (مقدار ثابت) ۴ است.
مجموع جملات تصاعد حسابی با فرمول زیر محاسبه می‌شود:
\[\underbrace{a + \cdots + b}_{n \,\, \textrm{عدد}} = \frac{n(a+b)}{2}\]
که $a$ عدد اول،
$b$ عدد آخر و
$n$ تعداد اعداد است.
برای مثال،
\[3+7+11+15=\frac{4 \cdot (3+15)}{2} = 36.\]
اساس فرمول این است: مجموع از $n$ عدد تشکیل شده و میانگین مقدار هر عدد $(a+b)/2$ است.

\index{تصاعد هندسی}\index{geometric progression}
یک \key{تصاعد هندسی} دنباله‌ای از اعداد است
که نسبت هر دو عدد متوالی در آن ثابت است.
برای مثال،
\[3,6,12,24\]
یک تصاعد هندسی با قدر نسبت ۲ است.
مجموع جملات تصاعد هندسی با فرمول زیر محاسبه می‌شود:
\[a + ak + ak^2 + \cdots + b = \frac{bk-a}{k-1}\]
که $a$ عدد اول،
$b$ عدد آخر و
نسبت اعداد متوالی $k$ است.
برای مثال،
\[3+6+12+24=\frac{24 \cdot 2 - 3}{2-1} = 45.\]

اثبات فرمول به این صورت است. فرض کنید:
\[ S = a + ak + ak^2 + \cdots + b .\]
با ضرب دو طرف در $k$:
\[ kS = ak + ak^2 + ak^3 + \cdots + bk,\]
حل معادله
\[ kS-S = bk-a\]
فرمول را به دست می‌دهد.

حالت خاص مجموع تصاعد هندسی فرمول زیر است:
\[1+2+4+8+\ldots+2^{n-1}=2^n-1.\]

\index{مجموع هارمونیک}\index{harmonic sum}

یک \key{مجموع هارمونیک} مجموعی به فرم زیر است:
\[ \sum_{x=1}^n \frac{1}{x} = 1+\frac{1}{2}+\frac{1}{3}+\ldots+\frac{1}{n}.\]

یک کران بالا برای مجموع هارمونیک برابر $\log_2(n)+1$ است.
به طور مشخص، می‌توانیم هر جمله $1/k$ را طوری تغییر دهیم که $k$ تبدیل به نزدیک‌ترین توان دو شود که از $k$ بزرگ‌تر نیست.
برای نمونه، وقتی $n=6$ است، می‌توان مجموع را به صورت زیر تخمین زد:
\[ 1+\frac{1}{2}+\frac{1}{3}+\frac{1}{4}+\frac{1}{5}+\frac{1}{6} \le
1+\frac{1}{2}+\frac{1}{2}+\frac{1}{4}+\frac{1}{4}+\frac{1}{4}.\]
این کران بالا از $\log_2(n)+1$ بخش تشکیل می‌شود
($1$، $2 \cdot 1/2$، $4 \cdot 1/4$ و غیره)،
و مقدار هر بخش حداکثر ۱ است.

\subsubsection{نظریه مجموعه‌ها}

\index{set theory@نظریه مجموعه‌ها}
\index{set@مجموعه}
\index{intersection@اشتراک}
\index{union@اجتماع}
\index{difference@تفاضل}
\index{subset@زیرمجموعه}
\index{universal set@مجموعه مرجع}
\index{complement@متمم}

یک \key{مجموعه} گردایه‌ای از عناصر است.
برای نمونه، مجموعه
\[X=\{2,4,7\}\]
شامل عناصر ۲، ۴ و ۷ است.
نماد $\emptyset$ نشان‌دهنده مجموعه تهی است،
و $|S|$ اندازه یک مجموعه $S$ را نشان می‌دهد،
یعنی تعداد عناصر موجود در مجموعه.
برای نمونه، در مجموعه بالا، $|X|=3$.

اگر یک مجموعه $S$ شامل عنصری به نام $x$ باشد،
می‌نویسیم $x \in S$،
و در غیر این صورت می‌نویسیم $x \notin S$.
برای نمونه، در مجموعه بالا
\[4 \in X \hspace{10px}\textrm{و}\hspace{10px} 5 \notin X.\]

\begin{samepage}
مجموعه‌های جدید را می‌توان با استفاده از عملگرهای مجموعه‌ای ساخت:
\begin{itemize}
\item \key{اشتراک} $A \cap B$ شامل عناصری است
که هم در $A$ و هم در $B$ هستند.
برای نمونه، اگر $A=\{1,2,5\}$ و $B=\{2,4\}$،
آنگاه $A \cap B = \{2\}$.
\item \key{اجتماع} $A \cup B$ شامل عناصری است
که در $A$ یا $B$ یا هر دو هستند.
برای نمونه، اگر $A=\{3,7\}$ و $B=\{2,3,8\}$،
آنگاه $A \cup B = \{2,3,7,8\}$.
\item \key{متمم} $\bar A$ شامل عناصری است
که در $A$ نیستند.
تفسیر متمم بستگی به \key{مجموعه مرجع} دارد، که شامل تمام عناصر ممکن است.
برای نمونه، اگر $A=\{1,2,5,7\}$ و مجموعه مرجع
$\{1,2,\ldots,10\}$ باشد، آنگاه $\bar A = \{3,4,6,8,9,10\}$.
\item \key{تفاضل} $A \setminus B = A \cap \bar B$
شامل عناصری است که در $A$ هستند ولی در $B$ نیستند.
دقت کنید که $B$ می‌تواند شامل عناصری باشد که در $A$ نیستند.
برای نمونه، اگر $A=\{2,3,7,8\}$ و $B=\{3,5,8\}$،
آنگاه $A \setminus B = \{2,7\}$.
\end{itemize}
\end{samepage}

اگر هر عنصر از $A$ به $S$ نیز تعلق داشته باشد،
می‌گوییم $A$ یک \key{زیرمجموعه} از $S$ است،
و با $A \subset S$ نشان داده می‌شود.
یک مجموعه $S$ همواره شامل $2^{|S|}$ زیرمجموعه است،
شامل مجموعه تهی.
برای نمونه، زیرمجموعه‌های مجموعه $\{2,4,7\}$ عبارتند از
\begin{center}
$\emptyset$،
$\{2\}$، $\{4\}$، $\{7\}$، $\{2,4\}$، $\{2,7\}$، $\{4,7\}$ و $\{2,4,7\}$.
\end{center}

برخی از مجموعه‌های پرکاربرد عبارتند از
$\mathbb{N}$ (اعداد طبیعی)،
$\mathbb{Z}$ (اعداد صحیح)،
$\mathbb{Q}$ (اعداد گویا) و
$\mathbb{R}$ (اعداد حقیقی).
مجموعه $\mathbb{N}$
بسته به موقعیت می‌تواند به دو صورت تعریف شود:
یا $\mathbb{N}=\{0,1,2,\ldots\}$
یا $\mathbb{N}=\{1,2,3,...\}$.

همچنین می‌توان مجموعه‌ای را با استفاده از قاعده‌ای به شکل
\[\{f(n) : n \in S\},\]
ساخت که در آن $f(n)$ یک تابع است.
این مجموعه شامل تمام عناصری به شکل $f(n)$ است
که $n$ عضوی از $S$ باشد.
برای مثال، مجموعه
\[X=\{2n : n \in \mathbb{Z}\}\]
شامل تمام اعداد صحیح زوج است.

\subsubsection{منطق}

\index{logic}
\index{negation}
\index{conjuction}
\index{disjunction}
\index{implication}
\index{equivalence}

مقدار یک عبارت منطقی یا
\key{درست} (1) یا \key{نادرست} (0) است.
مهم‌ترین عملگرهای منطقی عبارت‌اند از
$\lnot$ (\key{نقیض})،
$\land$ (\key{عطف})،
$\lor$ (\key{فصل})،
$\Rightarrow$ (\key{استلزام}) و
$\Leftrightarrow$ (\key{هم‌ارزی}).
جدول زیر معنای این عملگرها را نشان می‌دهد:

\begin{center}
\begin{tabular}{rr|rrrrrrr}
$A$ & $B$ & $\lnot A$ & $\lnot B$ & $A \land B$ & $A \lor B$ & $A \Rightarrow B$ & $A \Leftrightarrow B$ \\
\hline
0 & 0 & 1 & 1 & 0 & 0 & 1 & 1 \\
0 & 1 & 1 & 0 & 0 & 1 & 1 & 0 \\
1 & 0 & 0 & 1 & 0 & 1 & 0 & 0 \\
1 & 1 & 0 & 0 & 1 & 1 & 1 & 1 \\
\end{tabular}
\end{center}

عبارت $\lnot A$ مقداری مخالف با $A$ دارد.
عبارت $A \land B$ درست است اگر هر دو $A$ و $B$
درست باشند،
و عبارت $A \lor B$ درست است اگر $A$ یا $B$ یا هر دو
درست باشند.
عبارت $A \Rightarrow B$ درست است
هرگاه هر زمان که $A$ درست باشد، $B$ نیز درست باشد.
عبارت $A \Leftrightarrow B$ درست است
اگر $A$ و $B$ هر دو درست یا هر دو نادرست باشند.

\index{predicate}

یک \key{محمول} عبارتی است که بسته به پارامترهای خود
درست یا نادرست است.
محمول‌ها معمولاً با حروف بزرگ نشان داده می‌شوند.
به عنوان مثال، می‌توان محمول $P(x)$ را تعریف کرد
که دقیقاً زمانی درست است که $x$ یک عدد اول باشد.
با استفاده از این تعریف، $P(7)$ درست است ولی $P(8)$ نادرست است.

\index{quantifier}

یک \key{سور} یک عبارت منطقی را
به عناصر یک مجموعه مرتبط می‌کند.
مهم‌ترین سورها
$\forall$ (\key{به ازای هر}) و $\exists$ (\key{وجود دارد}) هستند.
برای مثال،
\[\forall x (\exists y (y < x))\]
بدین معناست که به ازای هر عنصر $x$ در مجموعه،
عنصری مانند $y$ در مجموعه وجود دارد
به‌طوری‌که $y$ کوچک‌تر از $x$ است.
این گزاره در مجموعه اعداد صحیح درست است،
اما در مجموعه اعداد طبیعی نادرست است.

با استفاده از نمادگذاری شرح داده شده،
می‌توان انواع گزاره‌های منطقی را بیان کرد.
برای مثال،
\[\forall x ((x>1 \land \lnot P(x)) \Rightarrow (\exists a (\exists b (a > 1 \land b > 1 \land x = ab))))\]
بدین معناست که اگر عدد $x$ بزرگ‌تر از 1
و غیر اول باشد،
آنگاه اعدادی مانند $a$ و $b$ وجود دارند
که بزرگ‌تر از $1$ هستند و حاصل‌ضرب آن‌ها برابر $x$ است.
این گزاره در مجموعه اعداد صحیح درست است.

\subsubsection{توابع}

تابع $\lfloor x \rfloor$ عدد $x$ را
به سمت پایین به یک عدد صحیح گرد می‌کند، و تابع
$\lceil x \rceil$ عدد $x$ را
به سمت بالا به یک عدد صحیح گرد می‌کند. برای مثال،
\[ \lfloor 3/2 \rfloor = 1 \hspace{10px} \textrm{و} \hspace{10px} \lceil 3/2 \rceil = 2.\]

توابع $\min(x_1,x_2,\ldots,x_n)$
و $\max(x_1,x_2,\ldots,x_n)$
کوچک‌ترین و بزرگ‌ترین مقادیر میان
$x_1,x_2,\ldots,x_n$ را به دست می‌دهند.
برای نمونه،
\[ \min(1,2,3)=1 \hspace{10px} \textrm{و} \hspace{10px} \max(1,2,3)=3.\]

\index{فاکتوریل}

\key{فاکتوریل} $n!$ می‌تواند به این صورت تعریف شود:
\[\prod_{x=1}^n x = 1 \cdot 2 \cdot 3 \cdot \ldots \cdot n\]
یا به صورت بازگشتی:
\[
\begin{array}{lcl}
0! & = & 1 \\
n! & = & n \cdot (n-1)! \\
\end{array}
\]

\index{عدد فیبوناچی}

\key{اعداد فیبوناچی}
%\footnote{فیبوناچی (حدود ۱۱۷۵--۱۲۵۰) یک ریاضیدان ایتالیایی بود.}
در موقعیت‌های بسیاری پدیدار می‌شوند.
آن‌ها می‌توانند به صورت بازگشتی بدین شکل تعریف شوند:
\[
\begin{array}{lcl}
f(0) & = & 0 \\
f(1) & = & 1 \\
f(n) & = & f(n-1)+f(n-2) \\
\end{array}
\]
نخستین اعداد فیبوناچی عبارتند از:
\[0, 1, 1, 2, 3, 5, 8, 13, 21, 34, 55, \ldots\]
همچنین یک فرمول بسته
برای محاسبه اعداد فیبوناچی وجود دارد که گاهی به آن
\index{فرمول بینه} \key{فرمول بینه} می‌گویند:
\[f(n)=\frac{(1 + \sqrt{5})^n - (1-\sqrt{5})^n}{2^n \sqrt{5}}.\]

\subsubsection{لگاریتم‌ها}

\index{لگاریتم}

\key{لگاریتم} یک عدد $x$
به صورت $\log_k(x)$ نشان داده می‌شود که در آن $k$ پایه
لگاریتم است.
طبق تعریف،
$\log_k(x)=a$ دقیقاً زمانی برقرار است که $k^a=x$ باشد.

یک ویژگی کاربردی لگاریتم‌ها این است
که $\log_k(x)$ برابر است با تعداد دفعاتی که
باید $x$ را بر $k$ تقسیم کنیم تا به
عدد ۱ برسیم.
برای نمونه، $\log_2(32)=5$
زیرا ۵ بار تقسیم بر ۲ نیاز است:

\[32 \rightarrow 16 \rightarrow 8 \rightarrow 4 \rightarrow 2 \rightarrow 1 \]

لگاریتم‌ها اغلب در تحلیل
الگوریتم‌ها به کار می‌روند، زیرا بسیاری از الگوریتم‌های کارآمد
در هر مرحله مقداری را نصف می‌کنند.
بنابراین، می‌توانیم کارایی این الگوریتم‌ها را
با استفاده از لگاریتم تخمین بزنیم.

لگاریتم حاصل‌ضرب عبارت است از:
\[\log_k(ab) = \log_k(a)+\log_k(b),\]
و در نتیجه،
\[\log_k(x^n) = n \cdot \log_k(x).\]
علاوه بر این، لگاریتم خارج‌قسمت برابر است با:
\[\log_k\Big(\frac{a}{b}\Big) = \log_k(a)-\log_k(b).\]
فرمول کاربردی دیگر به این صورت است:
\[\log_u(x) = \frac{\log_k(x)}{\log_k(u)},\]
و با استفاده از این، محاسبه
لگاریتم در هر پایه‌ای امکان‌پذیر است اگر روشی برای
محاسبه لگاریتم در یک پایه ثابت وجود داشته باشد.

\index{لگاریتم طبیعی}

\key{لگاریتم طبیعی} $\ln(x)$ یک عدد $x$
لگاریتمی است که پایه آن $e \approx 2.71828$ می‌باشد.
ویژگی دیگر لگاریتم‌ها این است که
تعداد ارقام یک عدد صحیح $x$ در مبنای $b$ برابر است با
$\lfloor \log_b(x)+1 \rfloor$.
برای نمونه، نمایش
$123$ در مبنای $2$ برابر با 1111011 است و
$\lfloor \log_2(123)+1 \rfloor = 7$.

\section{مسابقات و منابع}

\subsubsection{IOI}

المپیاد جهانی کامپیوتر (IOI)
یک مسابقه سالانه برنامه‌نویسی برای
دانش‌آموزان دوره متوسطه است.
هر کشور مجاز به اعزام تیمی متشکل از
چهار دانش‌آموز به مسابقه است.
معمولاً حدود ۳۰۰ شرکت‌کننده
از ۸۰ کشور حضور دارند.

مسابقه IOI شامل دو آزمون پنج‌ساعته است.
در هر دو آزمون، شرکت‌کنندگان باید سه مسئله الگوریتمی با دشواری‌های مختلف را حل کنند.
مسائل به زیرمسئله تقسیم می‌شوند و هر زیرمسئله نمره مشخص دارد.
حتی اگر شرکت‌کنندگان در قالب تیم اعزام شوند، رقابت انفرادی است.

سرفصل‌های IOI \cite{iois} مباحث مجاز در آزمون را تعیین می‌کند.
تقریباً تمام مباحث سرفصل IOI در این کتاب پوشش داده شده است.

شرکت‌کنندگان IOI از طریق مسابقات ملی انتخاب می‌شوند.
پیش از IOI، مسابقات منطقه‌ای متعددی برگزار می‌شود،
مانند المپیاد انفورماتیک بالتیک (BOI)،
المپیاد انفورماتیک اروپای مرکزی (CEOI)
و المپیاد انفورماتیک آسیا-اقیانوسیه (APIO).

برخی کشورها مسابقات تمرینی برخط برای شرکت‌کنندگان آینده IOI برگزار می‌کنند،
مانند مسابقات آزاد انفورماتیک کرواسی \cite{coci}
و المپیاد کامپیوتر آمریکا \cite{usaco}.
همچنین، مجموعه بزرگی از مسائل مسابقات لهستان به‌صورت برخط در دسترس است \cite{main}.

\subsubsection{ICPC}

مسابقه بین‌المللی برنامه‌نویسی دانشجویی (ICPC)
یک مسابقه سالانه برنامه‌نویسی برای دانشجویان است.
هر تیم شامل سه دانشجو است.
برخلاف IOI، اعضا با یکدیگر همکاری می‌کنند؛
تنها یک رایانه برای هر تیم وجود دارد.

مسابقه ICPC شامل چند مرحله است و در نهایت بهترین تیم‌ها به مرحله نهایی جهانی دعوت می‌شوند.
با وجود ده‌ها هزار شرکت‌کننده، تنها تعداد اندکی سهمیه نهایی\footnote{تعداد دقیق سهمیه‌های نهایی هر سال متغیر است؛ سال ۲۰۱۷، تعداد ۱۳۳ سهمیه نهایی وجود داشت.} وجود دارد،
بنابراین صعود به مرحله نهایی در برخی منطقه‌ها دستاوردی بزرگ محسوب می‌شود.

در هر مسابقه ICPC، تیم‌ها پنج ساعت زمان دارند تا حدود ده مسئله الگوریتمی را حل کنند.
راه‌حل مسئله تنها در صورتی پذیرفته می‌شود که همه آزمایه‌ها را به‌طور بهینه حل کند.
در طول مسابقه، رقبا می‌توانند نتایج تیم‌های دیگر را ببینند،
اما در ساعت پایانی جدول امتیازات منجمد می‌شود و دیدن نتایج آخرین ارسال‌ها ممکن نیست.

مباحث ممکن در ICPC به اندازه IOI مشخص نیست.
در هر صورت، برای ICPC به دانش بیشتری نیاز است، به‌ویژه مهارت‌های ریاضیاتی قوی‌تر.

\subsubsection{مسابقات برخط}

مسابقات برخط متعددی نیز برای همگان در دسترس است.
در حال حاضر فعال‌ترین پایگاه مسابقات Codeforces است
که تقریباً هفتگی مسابقه برگزار می‌کند.
در Codeforces، شرکت‌کنندگان دو بخش هستند:
تازه‌کارها در Div2 و برنامه‌نویسان با‌تجربه‌تر در Div1 رقابت می‌کنند.
سایر پایگاه‌های مسابقه عبارتند از: AtCoder، CS Academy، HackerRank و Topcoder.

برخی شرکت‌ها مسابقات برخط با مرحله نهایی حضوری برگزار می‌کنند.
نمونه‌های آن Facebook Hacker Cup،
Google Code Jam و Yandex.Algorithm است.
البته شرکت‌ها از این مسابقات برای استخدام نیرو نیز استفاده می‌کنند؛
عملکرد خوب در مسابقه روش مناسبی برای اثبات مهارت‌هاست.

\subsubsection{کتاب‌ها}

پیش از این نیز کتاب‌هایی (به‌جز این کتاب) با تمرکز بر برنامه‌نویسی رقابتی و حل مسئله الگوریتمی تالیف شده‌اند:

\begin{itemize}
\item S. S. Skiena and M. A. Revilla:
\emph{Programming Challenges: The Programming Contest Training Manual} \cite{ski03}
\item S. Halim and F. Halim:
\emph{Competitive Programming 3: The New Lower Bound of Programming Contests} \cite{hal13}
\item K. Diks et al.: \emph{Looking for a Challenge? The Ultimate Problem Set from
the University of Warsaw Programming Competitions} \cite{dik12}
\end{itemize}

دو کتاب نخست برای افراد مبتدی در نظر گرفته شده‌اند، در حالی که کتاب آخر شامل مباحث پیشرفته است.

البته کتاب‌های عمومی الگوریتم نیز برای برنامه‌نویسان رقابتی مناسب هستند.
برخی از کتاب‌های محبوب عبارتند از:

\begin{itemize}
\item T. H. Cormen, C. E. Leiserson, R. L. Rivest and C. Stein:
\emph{Introduction to Algorithms} \cite{cor09}
\item J. Kleinberg and É. Tardos:
\emph{Algorithm Design} \cite{kle05}
\item S. S. Skiena:
\emph{The Algorithm Design Manual} \cite{ski08}
\end{itemize}